%% Annexe A : Glossaire
\chapter{Glossaire}
\label{ann:glossaire}

Ce glossaire reprend les notions définies dans les chapitres. Le renvoi
indique le chapitre où la notion est présentée.

\begin{description}
    \item[Acquisition à chaud] Capture de données sur un système en fonctionnement, avant son extinction (Ch.~\ref{chap:acquisition}).
    \item[Bloqueur en écriture] Dispositif matériel ou logiciel qui autorise la lecture d'un support mais interdit toute modification (Ch.~\ref{chap:preserver}).
    \item[Chaîne de possession] Continuité documentée de la détention d'une pièce, de la saisie à la présentation (Ch.~\ref{chap:custody}).
    \item[Condensat] Empreinte de longueur fixe produite par une fonction de hachage, qui prouve l'intégrité (Ch.~\ref{chap:duplication}).
    \item[Effacement sécurisé] Réécriture de l'espace pour empêcher la récupération, laissant néanmoins des indices (Ch.~\ref{chap:anti}).
    \item[Fonction de hachage] Fonction transformant une donnée en un condensat, sensible au moindre bit (MD5, SHA-256) (Ch.~\ref{chap:duplication}).
    \item[Forensique numérique] Investigation numérique à vocation probatoire, soumise à la rigueur de la preuve (Ch.~\ref{chap:forensique}).
    \item[Image forensique] Copie bit à bit d'un support, reproduisant fichiers, espace non alloué et métadonnées (Ch.~\ref{chap:duplication}).
    \item[Impartialité] Disposition à suivre les preuves où elles mènent, sans favoriser une thèse (Ch.~\ref{chap:ethique}).
    \item[Investigation numérique] Démarche générale de collecte, d'analyse et d'interprétation d'éléments numériques (Ch.~\ref{chap:forensique}).
    \item[Métadonnées] Données portées par un fichier (dates, nom, propriétaire), utiles pour dater une action (Ch.~\ref{chap:disque}).
    \item[Neutralité] Absence d'intérêt personnel de l'enquêteur dans l'issue (Ch.~\ref{chap:ethique}).
    \item[Récupération par signatures] Reconstitution de fichiers par leurs motifs de début et de fin, sans l'index (Ch.~\ref{chap:artefacts}).
    \item[Recevabilité] Aptitude d'une preuve à être admise par la justice, selon la manière dont elle a été obtenue et préservée (Ch.~\ref{chap:recevabilite}).
    \item[Signature (de fichier)] Octets de début d'un fichier qui trahissent son type, indépendamment de l'extension (Ch.~\ref{chap:artefacts}).
    \item[Système de fichiers] Organisation des données d'une partition (FAT, NTFS, EXT) (Ch.~\ref{chap:disque}).
    \item[Trilemme CRO] Confidentialité, Fiabilité, Opposabilité~: les trois dimensions de toute trace et de tout acte d'enquête (Ch.~\ref{chap:prologue}).
    \item[Verdict à trois valeurs] OUI (établi), NON (réfuté), NEK (indéterminé) (Ch.~\ref{chap:mesurer}).
\end{description}
